% =====================================================================
% From Latent-Diffusion to Patch-Diffusion: a failure analysis and redesign
% Standalone, PDF-ready. Standard packages only; no \includegraphics, so it
% compiles anywhere. Figures referenced by filename live in:
%   brats-uad/report/figures/   (latent approach)
%   04_PDM/report_pdm/figures/  (patch approach)
%
% Compile:  pdflatex latent_to_patch.tex   (run twice for the ToC)
%        or upload to Overleaf.
%
% Every number here is taken from the actual code / logs:
%   brats-uad/  (latent, "Gen 2")  -> local_logs/diffusion_*/run.log,
%               config_snapshot.json, report/main.tex
%   04_PDM/     (patch,  "Gen 3")  -> outputs/logs/*, calibration.json, config.py
% =====================================================================
\documentclass[11pt,a4paper]{article}

\usepackage[utf8]{inputenc}
\usepackage[a4paper,margin=22mm]{geometry}
\usepackage{amsmath,amssymb,mathtools}
\usepackage{bm}
\usepackage{booktabs}
\usepackage{array}
\usepackage{longtable}
\usepackage{xcolor}
\usepackage{enumitem}
\usepackage{titlesec}
\usepackage[colorlinks=true,linkcolor=blue!50!black,urlcolor=red!50!black]{hyperref}

\newcommand{\x}{\bm{x}}
\newcommand{\z}{\bm{z}}
\newcommand{\R}{\mathbb{R}}
\newcommand{\code}[1]{\texttt{\small #1}}
\newcommand{\abar}{\bar{\alpha}}

\title{\textbf{From Latent-Diffusion to Patch-Diffusion}\\[3pt]
\large Why the latent approach stalled, and what the patch redesign fixes}
\author{Evidence-grounded analysis of the \code{brats-uad} (latent) and
\code{04\_PDM} (patch) runs}
\date{\today}

\begin{document}
\maketitle

\begin{abstract}
\noindent
This document explains, with the real numbers from both runs, why the
\emph{latent-diffusion} approach to unsupervised anomaly detection (UAD) on
BraTS-PEDs reached a hard ceiling, and how the \emph{patch-based pixel-diffusion}
redesign removes the specific bottleneck that caused it. The latent system
detected lesions well (slice AUROC $\approx0.86$--$0.94$) but produced an almost
unusable segmentation (calibrated Dice $\approx0.00$--$0.10$). This document traces this
\emph{AUROC--Dice gap} to a concrete, architectural cause --- a brain-boundary
``edge rim'' generated by the $\times8$ KL-VAE --- and show that the patch
approach eliminates that cause at its root by abandoning the VAE entirely and
working in full-resolution pixel space on overlapping patches. The honest
conclusion is not ``patch scores higher'' (the evaluation sets differ); it is
that the limiting factor moves from a fixed \emph{architectural artefact} to a
\emph{data-coverage} problem that more training data can solve.
\end{abstract}

\tableofcontents
\clearpage

% =====================================================================
\section{Context: three generations on the same problem}
% =====================================================================
The task is fixed throughout: train a generative model on \emph{healthy}
pediatric brain MRI only, then flag tumours in unseen patients as deviations
from the learned healthy manifold --- no lesion labels in training. The
implementation went through three generations:

\begin{longtable}{@{}rlll@{}}
\toprule
Gen & Codebase & Approach & Outcome \\
\midrule
\endhead
1 & \code{01\_cv-project} & Latent diffusion on Stable Diffusion's RGB VAE & abandoned: natural-image bias, dropped $t1n$, $\pm3\sigma$ clip \\
2 & \code{brats-uad}      & Latent diffusion on a from-scratch 4-ch medical KL-VAE & \textbf{the latent result analysed here} \\
3 & \code{04\_PDM}        & Patch-based pixel diffusion (no VAE) & \textbf{the current approach} \\
\bottomrule
\end{longtable}

\noindent
Generation~2 already fixed the obvious Generation~1 mistakes (it built a
medical 4-channel VAE, kept all four modalities plus a contrast-enhancement
channel, and used robust percentile normalisation). It is therefore the
\emph{strongest} version of the latent idea, and the right baseline to learn
from. This document compares Generation~2 (latent) against Generation~3 (patch).

% =====================================================================
\section{The latent-diffusion design (Generation 2)}
% =====================================================================

\subsection{Pipeline}
\[
\underbrace{\x\in\R^{4\times256\times256}}_{\text{slice}}
\xrightarrow{\;E_\phi\;}
\underbrace{\z\in\R^{4\times32\times32}}_{\text{latent}}
\xrightarrow[\text{DDIM}]{\text{noise to }T_{\text{int}},\ \text{denoise}}
\hat\z
\xrightarrow{\;D_\psi\;}
\hat\x
\;\Rightarrow\;
\text{residual }|\x-\hat\x|.
\]

\subsection{The medical KL-VAE}
A from-scratch four-channel KL-regularised VAE
(\code{brats-uad/src/models/kl\_vae.py}): GroupNorm--SiLU ResNet blocks with a
single self-attention block at the bottleneck, downsampling by a factor of
\textbf{eight}, $\x\in\R^{4\times256\times256}\to\z\in\R^{4\times32\times32}$.
Key settings (from the report and \code{config\_snapshot.json}):
\begin{itemize}[leftmargin=1.4em,itemsep=0pt]
\item Reconstruction loss $\mathcal L_{\text{rec}}=\|\x-\hat\x\|_1+0.5\,(1-\text{MS-SSIM}(\x,\hat\x))$;
\item KL weight $\beta_{\text{KL}}=10^{-6}$ (deliberately tiny --- the bottleneck,
not the prior, does the regularising, so the latent keeps reconstruction detail);
\item Empirical latent scaling $s=1/\sigma_z=1.262$ (measured $\sigma_z=0.793$ over
$2{,}000$ slices), replacing Stable Diffusion's $0.18215$.
\end{itemize}

\subsection{Latent DDPM and dual-space residual}
A DDPM is trained on the \emph{standardised healthy latents} ($T=1000$, cosine
$\beta\in[10^{-4},2\times10^{-2}]$, batch $32$, lr $10^{-4}$, EMA $0.999$). At
inference, partial-noise DDIM ($50$ steps) denoises the latent back to the
healthy manifold; the decoded reconstruction is the pseudo-healthy
counterfactual. The score fuses a \emph{pixel} residual with a \emph{latent}
residual (channel-averaged latent distance, upsampled), weight $\alpha=0.5$ ---
the latent term was added precisely because the decoder can hallucinate healthy
texture and mask a pixel-space anomaly.

% =====================================================================
\section{What the latent approach achieved}
% =====================================================================

\subsection{Training behaved correctly}
From \code{local\_logs/diffusion\_20260618\_134532/run.log}: the noise-prediction
MSE fell from $0.61$ to $0.45$ on train; the validation denoising loss bottomed
at \textbf{$0.411$ at epoch 8}, and early stopping halted the run at epoch 18
(patience 10). The validation loss sat \emph{below} the training loss throughout
(overfit gap $-4\%$ to $-20\%$). \textbf{Reading:} the model is in the
under-/well-fit regime, not memorising the pilot set --- exactly what a
healthy-manifold model should do. So whatever goes wrong downstream is
\emph{not} a training pathology.

\subsection{Detection worked; segmentation did not}
Per-slice metrics on test patient \code{BraTS-PED-00029} (pilot):

\begin{longtable}{@{}lccl@{}}
\toprule
Axial slice & AUROC & Dice (calibrated) & Lesion \\
\midrule
\endhead
$z=065$ & 0.863 & \textbf{0.000} & yes (small) \\
$z=071$ & 0.922 & 0.046 & yes \\
$z=075$ & 0.927 & 0.077 & yes \\
$z=080$ & 0.944 & 0.099 & yes \\
\midrule
\textbf{mean} & \textbf{$\approx$0.91} & \textbf{$\approx$0.06} & --- \\
\bottomrule
\end{longtable}

\noindent
A chance detector scores AUROC $0.50$; $0.91$ is strong slice-level
discrimination. Yet the deployable Dice is near zero. That divergence is the
whole story.

% =====================================================================
\section{The core problem: the AUROC--Dice gap}
% =====================================================================

\subsection{What the gap rules out}
A high AUROC with a low Dice is diagnostic. It rules out the two trivial failure
modes immediately:
\begin{itemize}[leftmargin=1.4em,itemsep=0pt]
\item \textbf{Not a detection failure}: AUROC $\approx0.91$ means the continuous
score \emph{ranks} tumour voxels above healthy voxels --- the lesion signal is
present and correctly ordered.
\item \textbf{Not mis-training}: the loss curves are clean and non-overfit
(\S3.1).
\end{itemize}
So the lesion \emph{is} found by the score, but the \emph{binary} prediction at a
single global threshold is wrong.

\subsection{The mechanism: two super-threshold structures}
The anomaly residual contains \emph{two} regions that both exceed the
healthy-calibrated threshold:
\begin{enumerate}[leftmargin=1.5em,itemsep=0pt]
\item the \textbf{tumour} (the structure of interest), and
\item a \textbf{brain-boundary rim} (an unwanted artefact).
\end{enumerate}
A single global $95$th-percentile threshold cannot separate them. Because the rim
encircles the whole brain, thresholding \textbf{over-selects the boundary and
under-selects the compact tumour}, so $|\hat P\cap G|$ stays small while
$|\hat P|$ is large --- and Dice $=2|\hat P\cap G|/(|\hat P|+|G|)$ collapses
toward zero. The hot-spot is visually correct (which is why AUROC is high), but
the binary mask is dominated by the rim.

\subsection{Where the rim comes from: the VAE, not the diffusion model}
This is the key finding. The rim is \emph{not} produced by the diffusion process;
it is produced by the autoencoder. An $\times8$ VAE systematically
\textbf{under-reconstructs the high-contrast tissue/background and skull
interface}, because that sharp edge is exactly the high-frequency content a
lossy $8\times$ bottleneck discards. The result is a large but
\emph{anatomically constant} reconstruction error all along the brain boundary
(visible in \code{brats-uad/report/figures/fig\_vae\_recon.png} and
\code{results/vae\_reconstruction\_analysis.png}). This edge energy is comparable
in magnitude to the tumour residual, so the healthy-percentile threshold passes
both. Erosion of the brain mask and a healthy baseline $M_k$ only partially
suppress it on a 30-patient calibration set.

\subsection{Why this is intrinsic, not a tuning bug}
Three independent properties of \emph{any} latent-VAE UAD make the rim hard to
remove by tuning:
\begin{enumerate}[leftmargin=1.5em,itemsep=1pt]
\item \textbf{Lossy bottleneck.} An $\times8$ compression \emph{must} discard
high-frequency detail; the sharpest edge in the image (brain/background) is the
first casualty, so a boundary residual is structurally guaranteed.
\item \textbf{Spatial-precision floor.} The anomaly map is born at $32\times32$
and upsampled; its boundary sharpness is capped regardless of the diffusion
model's quality. Small or thin lesions are below this floor.
\item \textbf{Decoder hallucination.} A decoder good enough to reconstruct
healthy tissue can also ``paint in'' plausible tissue over a real anomaly,
shrinking the very residual being measured (the reason a latent-space residual had to
be fused in as a patch).
\end{enumerate}
In short, the latent design couples detection (which works) to a reconstruction
artefact and a resolution floor (which cap the achievable Dice). No threshold can
undo it while the rim and the lesion coexist above threshold.

% =====================================================================
\section{The patch-based redesign (Generation 3) and why it fixes each problem}
% =====================================================================

\subsection{The one structural change}
Delete the VAE. Do diffusion directly on $96\times96$ \emph{pixel} patches
(\code{04\_PDM}). Everything else (healthy-only training, partial-noise DDIM
counterfactual, modality-weighted CE-aware residual, healthy-percentile
calibration) is preserved or improved. Removing the autoencoder removes the
\emph{source} of the rim rather than trying to subtract it afterwards.

\subsection{Problem-to-fix mapping}
\begin{longtable}{@{}p{0.46\textwidth}p{0.50\textwidth}@{}}
\toprule
Latent problem & Patch fix (\code{04\_PDM}) \\
\midrule
\endhead
VAE generates a boundary \textbf{edge rim} that dominates the residual &
\textbf{No VAE.} Residual is taken directly in pixel space $|\x-\hat\x|$; the rim
has no source. \\
Whole-slice encode/decode makes \emph{one} global boundary artefact &
Work on \textbf{$96\times96$ patches} (stride $16$, $121$/slice); fuse with a
\textbf{centre-weighted Gaussian} ($\sigma=32$) that down-weights each patch's
edges, so boundary voxels never dominate. \\
$\times8$ latent caps boundary sharpness &
\textbf{Full pixel resolution} --- the score is computed at native $256\times256$,
no upsampling floor. \\
Decoder can hallucinate healthy texture over a lesion &
No decoder; the denoiser only ever sees and edits the pixel patch directly. \\
DDPM trained on few whole-slice latents &
$121$ patches/slice turns $\sim$3k healthy slices into $\sim$570k (uncapped:
$\sim$2.09M) training samples --- \textbf{data-rich} in the 30-patient regime, and
forces a \emph{local} texture model that cannot memorise global layout. \\
Gaussian noise attacks pixel detail, not lesion scale &
\textbf{Simplex (correlated) noise} matched to lesion size, so lesions are
in-painted as healthy tissue and survive as a residual. \\
Single noise level &
\textbf{Multi-scale} scoring at $T\in\{50,150,250\}$ adapts sensitivity to lesion
size. \\
\bottomrule
\end{longtable}

\subsection{The patch model in numbers}
From \code{04\_PDM/outputs/}: a $25.3$\,M-parameter pixel U-Net (blocks
$64{-}128{-}256{-}256$, attention only at the $12\times12$ bottleneck), trained on
an A100-80GB, best validation noise-MSE \textbf{$0.0604$ at epoch 58},
calibrated threshold $\tau=0.128$. Over 204 test slices (4 patients) it reaches
AUROC $0.713$, Dice $0.076$ (oracle $0.138$); on its best-represented patient,
AUROC $0.965$ / Dice $0.589$.

\subsection{The bottleneck moved (the real point)}
The latent approach was limited by a \emph{fixed architectural artefact} (the
VAE rim) --- something no amount of extra data removes. The patch approach has no
such artefact; its failures are now \emph{out-of-manifold patients} whose healthy
anatomy the 30-patient training set never covered (the residual is diffuse over
the whole brain, not a boundary ring). That is a \textbf{data-coverage} problem,
which more training patients directly fix. The limiting factor moved from
``the architecture'' to ``the dataset size'' --- a more honest and more
addressable place to be.

% =====================================================================
\section{An honest caveat on comparing the two}
% =====================================================================
It would be wrong to claim ``patch beats latent'' from the headline AUROCs. The
latent figure ($\approx0.91$) is from a \emph{single, well-behaved test patient}
(\code{BraTS-PED-00029}), reported qualitatively. The patch figure ($0.713$) is a
\emph{four-patient average that deliberately includes a failure case}
(\code{P00216}, AUROC $0.571$). The two numbers are not measured on the same
population and must not be put side by side as a win.

The defensible comparison is \emph{qualitative and mechanistic}:
\begin{itemize}[leftmargin=1.4em,itemsep=0pt]
\item The latent approach's Dice was capped by a \textbf{VAE edge artefact} that a
single threshold could not separate from the lesion.
\item The patch approach \textbf{removes that artefact at its source}, so where
the score separates the classes the calibrated threshold is near-optimal (within
$0.005$ Dice of the sweep optimum), and the remaining gap is attributable to
training-set coverage, not to the detector.
\end{itemize}
Both remain 30-patient pilots; neither has a full-cohort, statistically powered
evaluation yet.

% =====================================================================
\section{Summary}
% =====================================================================
\begin{longtable}{@{}p{0.30\textwidth}p{0.32\textwidth}p{0.32\textwidth}@{}}
\toprule
& \textbf{Latent (Gen 2, \code{brats-uad})} & \textbf{Patch (Gen 3, \code{04\_PDM})} \\
\midrule
Representation     & $\times8$ KL-VAE latent ($4{\times}32{\times}32$) & raw pixels, $96{\times}96$ patches \\
Diffusion space    & latent manifold & pixel manifold \\
Noise              & Gaussian & simplex (correlated) \\
Score              & pixel + latent residual ($\alpha{=}0.5$) & multi-scale pixel residual + CE, Gaussian-fused \\
Best val noise-MSE & $0.411$ (latent space, ep 8) & $0.0604$ (pixel space, ep 58) \\
Detection (AUROC)  & $\approx0.91$ (1 patient) & $0.713$ (4 patients; $0.965$ best) \\
Calibrated Dice    & $\approx0.06$ & $0.076$ (best slice $0.589$) \\
Limiting factor    & \textbf{VAE edge-rim artefact} (architectural, fixed) & \textbf{healthy-manifold coverage} (data, fixable) \\
\bottomrule
\end{longtable}

\noindent
\textbf{Bottom line.} The latent-diffusion pilot proved the counterfactual
principle works (it detected lesions) but exposed a structural ceiling: an
$\times8$ VAE injects a brain-boundary residual that destroys segmentation
precision. The patch-diffusion redesign removes the VAE, computes the residual in
full-resolution pixel space on overlapping patches, and thereby eliminates that
ceiling --- shifting the remaining limitation onto training-set size, which the
full BraTS-PEDs cohort can address.
\end{document}
